\section{Data Preprocessing}
To construct an academically defensible and deployment-ready forecasting pipeline, we enforce strict data hygiene and temporal constraints to eliminate out-of-sample data leakage:

\begin{itemize}
    \item \textbf{Target Shift Formulation:} The forecast target \texttt{risk\_level} is shifted by $-1$ observation (a 2-hour look-ahead interval). All model input features are strictly restricted to time step $t$ to predict the grid risk level at step $t+1$ ($T+2\text{h}$).
    \item \textbf{Data Curation and Row Splitting:} Shifting the target variable by $-1$ per feeder stream generates exactly 1,000 NaN values (corresponding to the final timestamp of each of the 1,000 distinct feeder streams), which are dropped. This reduces the raw dataset from 65,983 to 64,983 valid telemetry records.
    \item \textbf{Chronological Train-Test Split:} Telemetry records are grouped by \texttt{substation\_id} and \texttt{feeder\_id} and sorted chronologically. For each feeder stream, the first 80\% of observations are assigned to the training split ($N = 51,587$), and the final 20\% are allocated to the hold-out test split ($N = 13,396$). Random shuffling and K-fold cross-validation are strictly disabled to preserve real-world sequence dynamics.
    \item \textbf{3D Sliding Window Generation:} Models ingest historical sequences of length $seq\_len = 5$ (5 steps $\times$ 2 hours = 10-hour lookback window). In the hold-out test set, the initial 4 observations per feeder stream cannot form a complete 5-step sequence and are omitted from sequence-based evaluation. This yields exactly 42,674 training, 10,545 validation, and 10,027 test sequence windows ($X \in \mathbb{R}^{10027 \times 5 \times 28}$).
    \item \textbf{Leakage-Free Feature Scaling:} Feature standardization via \texttt{StandardScaler} is fitted solely on the training partition. The derived mean and variance parameters are subsequently applied to scale the validation and test partitions, preventing forward temporal leakage.
\end{itemize}
